\documentclass[a4paper,12pt]{report}
\usepackage[utf8]{inputenc}
\usepackage[portuguese]{babel}
\usepackage[style=apa, backend=biber]{biblatex}
\addbibresource{individual.bib} % Nome do arquivo de bibliografia
\usepackage{csquotes} % Necessário para citações com babel
\usepackage{graphicx}
\usepackage{amsmath}
\usepackage{hyperref}
\usepackage{geometry}
\usepackage{setspace}
\usepackage{indentfirst} % Pacote para indentar o primeiro parágrafo
\usepackage{newtxtext,newtxmath} % Usa fontes Times modernas
\usepackage{xcolor} % Pacote para cores
\usepackage{fancyhdr} % Para personalização do rodapé
\usepackage{titlesec} % Para personalizar títulos e subtítulos

% Configuração de Margens e Espaçamento
\geometry{a4paper, margin=2.5cm}
\setstretch{1.5}

% Configuração de Tamanhos de Fonte e Espaçamento para Seções
\titleformat{\section} % Títulos principais numerados
{\normalfont\fontsize{14}{17}\bfseries}{\thesection}{1em}{}
\titleformat*{\section} % Títulos principais não numerados
{\normalfont\fontsize{14}{17}\bfseries}
\titleformat{\subsection} % Subtítulos
{\normalfont\fontsize{12}{15}\bfseries}{\thesubsection}{1em}{}

% Ajuste de Espaçamento para Seções
\titlespacing{\section}
{0pt}{1.5em}{1em}
\titlespacing*{\section}
{0pt}{1.5em}{1em}
\titlespacing{\subsection}
{0pt}{1.2em}{0.8em}

% Configuração de Rodapé com Numeração
\pagestyle{fancy}
\fancyhf{} % Limpa cabeçalho e rodapé
\fancyfoot[R]{\thepage} % Numeração no canto inferior direito
\renewcommand{\headrulewidth}{0pt} % Remove linha do cabeçalho
\renewcommand{\footrulewidth}{0pt} % Remove linha do rodapé

% Configuração de Cores
\definecolor{barraazul}{RGB}{0, 51, 153}

% Numeração das seções sem números iniciais adicionais
\renewcommand{\thesection}{\arabic{section}}
\renewcommand{\thesubsection}{\arabic{section}.\arabic{subsection}}
\renewcommand{\thesubsubsection}{\arabic{section}.\arabic{subsection}.\arabic{subsubsection}}

% Redefinição do Cabeçalho da Bibliografia
\defbibheading{bibliography}{
	\section*{Referências Bibliográficas}
	\addcontentsline{toc}{section}{Referências Bibliográficas}
	\thispagestyle{fancy} % Garante que o rodapé seja consistente
}

% Garante que o estilo de página seja 'fancy' na bibliografia
\AtBeginBibliography{\pagestyle{fancy}}

% Ajuste para indentação de parágrafos
\setlength{\parindent}{1.25cm}

\begin{document}
	
	% Capa
\begin{titlepage}
	\centering
	\vspace*{-2cm} % Reduz a margem superior da página de título
	
	% Ajusta as imagens para remover espaçamentos internos (se houver)
	\raisebox{-0.5\height}{\includegraphics[width=0.4\textwidth]{iscte.png}}%
	\hfill%
	\raisebox{-0.46\height}{\includegraphics[width=0.4\textwidth]{ista.png}}\\[0.5cm]
	
	\noindent
	{\color{barraazul}\rule{\textwidth}{1mm}} % Barra azul horizontal
	\\[1cm]
	
	{\LARGE  \textbf{A Influência das Emoções na Tomada de Decisão, Uso de Heurísticas e os Desafios na Inteligência Artificial Emocional} \par}
	\vspace{3cm}
	
	{\Large \textbf{Mestrado em Inteligência Artificial}} \par
	{\large \textbf{Unidade Curricular: Cognição e Emoção}} \par
	\vspace{3cm}
	
	{\large \textbf{Aluno: Luís Ricardo Silva Inácio}} \par
	{\large \textbf{Número: 129074}} \par
	\vspace{3cm}
	

	
	{\large \textbf{Professora: Doutora Cristiane da Anunciação Souza}} \par
	\vfill
	
	\vspace{0.5cm}
	{\large \textbf{Data de Entrega: 29 de novembro de 2024}} \par
\end{titlepage}

	
	% Costas da Capa (Em branco)
	\newpage
	\thispagestyle{empty}
	\mbox{}
	\newpage
	
	% Configuração para numeração romana
	\pagenumbering{roman}
	
	% Resumo
	\section*{Resumo}
	\addcontentsline{toc}{section}{Resumo}
	
	Este artigo analisa de forma abrangente o papel das emoções no uso de heurísticas e como a complexidade e ambiguidade da linguagem representam desafios na criação de sistemas de inteligência artificial (IA) emocional. Inicialmente, discute-se como as emoções influenciam as heurísticas cognitivas, facilitando decisões rápidas, mas também podendo levar a vieses e erros. Explora-se a dificuldade que sistemas de IA enfrentam na interpretação da linguagem emocional humana, devido a nuances como sarcasmo, ironia e expressões culturais. O trabalho examina tanto os impactos positivos das emoções na tomada de decisão e na IA—como a melhoria na interação homem-máquina—quanto os impactos negativos, incluindo exemplos práticos como assistentes virtuais e chatbots. Questões éticas e de privacidade são aprofundadas, destacando os riscos de manipulação emocional e invasão de privacidade na recolha de dados emocionais. Por fim, propõem-se abordagens para mitigar esses riscos, enfatizando a importância de modelos explicáveis, conformidade regulatória e educação dos utilizadores. Conclui-se que o desenvolvimento de IA emocional requer uma abordagem ética e interdisciplinar, equilibrando avanços tecnológicos com o respeito pelos valores humanos.
	
	\vspace{4em}
	
	\noindent\textbf{Palavras-Chave:} \normalsize{Emoções, Heurísticas Cognitivas, Tomada de Decisão, Inteligência Artificial Emocional, Privacidade de Dados, Processamento de Linguagem Natural, Ética em IA.}
	
	\newpage	
	
	% Configuração para numeração arábica
	\pagenumbering{arabic}
	
	% Introdução
	\section{Introdução}
	
	As emoções desempenham um papel fundamental na forma como os indivíduos processam informações e tomam decisões. Elas influenciam directamente o uso de heurísticas, estratégias cognitivas que facilitam a tomada de decisões em contextos complexos ou incertos \parencite{kahneman1974}. Paralelamente, a linguagem, como meio de expressão emocional, é inerentemente complexa e ambígua, representando um desafio significativo para a criação de sistemas de inteligência artificial (IA) capazes de compreender e responder às emoções humanas \parencite{chomsky1965}. Este artigo explora a intersecção entre emoções, heurísticas e linguagem na cognição humana e discute as implicações técnicas, éticas e de privacidade na criação de sistemas de IA emocional. Além disso, propõe abordagens para mitigar os riscos associados e promover o desenvolvimento de IA emocionalmente inteligente e ética, alinhada com as directrizes de design ético propostas pelo \textcite{ieee2019}.

	\section{O Papel das Emoções nas Heurísticas e a Complexidade da Linguagem na IA Emocional}
	
	\subsection{Heurísticas, Emoções e Linguagem: Uma Introdução Teórica}
	
	\textcite{kahneman1974} definiram heurísticas como mecanismos que permitem aos indivíduos lidar com problemas de forma eficiente, utilizando uma quantidade limitada de recursos cognitivos. Estas estratégias são particularmente úteis em situações em que o tempo ou as informações disponíveis são escassas. No entanto, as heurísticas não são infalíveis e podem levar a erros ou vieses cognitivos. As emoções influenciam directamente o uso de heurísticas, funcionando como atalhos que podem ajudar ou prejudicar decisões, dependendo do contexto \parencite{slovic2007}.
	
	A linguagem é o meio pelo qual as emoções são frequentemente expressas e interpretadas. Conforme descrito por \textcite{chomsky1965}, a linguagem humana é complexa e permite uma variedade infinita de expressões. Esta complexidade apresenta desafios significativos para a IA, que necessita de interpretar correctamente nuances culturais, contextuais e emocionais para compreender as emoções humanas e responder de forma adequada. Apesar dos avanços no processamento de linguagem natural, a IA ainda apresenta limitações para entender e replicar estas subtilezas, o que pode gerar lacunas e equívocos na interpretação emocional \parencite{russell2020}.
	
	\subsection{Impactos Positivos das Emoções no Uso de Heurísticas e IA Emocional}
	
	As emoções podem facilitar decisões rápidas e eficientes através de heurísticas. \textcite{isen2001} destaca que estados emocionais positivos, como a felicidade, podem ampliar a criatividade e a flexibilidade cognitiva. Indivíduos felizes tendem a adoptar uma abordagem mais global, resultando na aplicação de heurísticas adaptativas em contextos complexos. Em ambientes de trabalho, gestores com estados emocionais positivos são mais propensos a tomar decisões inovadoras e a resolver problemas de forma eficaz.
	
	Em sistemas de IA, a compreensão e incorporação de estados emocionais podem melhorar a interacção humano-máquina, tornando-a mais natural e eficaz. Assistentes virtuais que reconhecem e respondem a emoções proporcionam experiências de utilizador mais satisfatórias \parencite{picard1997}. Por exemplo, um sistema que detecta frustração no utilizador pode ajustar a sua abordagem para fornecer assistência de forma mais empática.
	
	\subsection{Desafios na Interpretação da Linguagem Emocional pela IA}
	
	A ambiguidade e a complexidade da linguagem emocional representam obstáculos significativos para a IA emocional. Frases como “Estou bem” podem variar em significado, desde sinceridade até sarcasmo, dependendo do tom, da entoação e do contexto. Modelos de processamento de linguagem natural actuais, embora tenham avançado na interpretação de texto, ainda enfrentam dificuldades na compreensão das nuances emocionais implícitas \parencite{russell2020}.
	
	A IA frequentemente interpreta a linguagem de forma literal, falhando em captar sarcasmo, ironia ou humor. As emoções influenciam directamente a forma como as pessoas escolhem palavras e estruturam frases, tornando a interpretação automática ainda mais desafiante. Por exemplo, uma mensagem irónica pode ser interpretada erroneamente como literal, levando a respostas inadequadas.
	
	\subsection{Impactos Negativos das Emoções e da Linguagem Ambígua na IA e na Tomada de Decisão}
	
	Embora as emoções possam facilitar decisões, também podem introduzir vulnerabilidades significativas. Emoções intensas podem distorcer julgamentos, levando a decisões enviesadas ou irracionais. Na IA, os sistemas que procuram simular ou interpretar emoções enfrentam o desafio de evitar a reprodução desses enviesamentos. \textcite{russell2020} alertam que modelos que imitam heurísticas emocionais humanas podem inadvertidamente incorporar vieses culturais ou sociais, prejudicando a fiabilidade e a ética das decisões automatizadas. Por exemplo, um \textit{chatbot} que aprende a partir de interacções humanas sem filtragem pode adoptar linguagem discriminatória ou ofensiva.
	
	A linguagem ambígua exacerba estes desafios. Sistemas de IA podem interpretar mal o sarcasmo, a ironia ou o humor, resultando em respostas inadequadas ou até prejudiciais. Esta falha na interpretação pode levar a decisões empresariais baseadas em dados imprecisos, afectando estratégias de negócio e satisfação do cliente.
	
	\subsection{Implicações Éticas e de Privacidade na Criação de IA Emocional}
	
	O desenvolvimento de sistemas de IA emocional levanta questões éticas significativas. A manipulação emocional é uma preocupação central; sistemas que monitorizam emoções podem explorar vulnerabilidades dos utilizadores, especialmente em contextos comerciais, influenciando decisões de compra e violando a autonomia dos consumidores. Além disso, a recolha de dados emocionais sensíveis levanta questões de privacidade. Há o perigo de que tais dados sejam utilizados para vigilância, discriminação ou outros fins maliciosos.
	
	A falta de transparência nos sistemas de IA emocional dificulta a compreensão de como as decisões são tomadas, comprometendo a confiança dos utilizadores \parencite{russell2020}. Sem uma explicação clara de como um sistema interpreta emoções e toma decisões, é difícil identificar e corrigir possíveis vieses ou erros.
	
	De acordo com as directrizes propostas pelo \textcite{ieee2019}, é fundamental que o desenvolvimento de sistemas autónomos e inteligentes priorize o bem-estar humano, incorporando princípios éticos desde a fase de concepção.
	
	\subsection{Abordagens para Mitigar Riscos Éticos e de Privacidade}
	
	Para abordar estes desafios, é essencial o desenvolvimento de modelos explicáveis que permitam aos utilizadores entender como as emoções são detectadas e influenciam as decisões do sistema \parencite{russell2020}. A explicabilidade aumenta a confiança do utilizador e permite a identificação de possíveis vieses ou falhas no sistema.
	
	A aplicação rigorosa de regulamentações como o Regulamento Geral de Protecção de Dados (RGPD) é fundamental para proteger a privacidade dos utilizadores. Técnicas como aprendizagem federada podem reduzir riscos de privacidade, permitindo treinar modelos sem transferir dados sensíveis para servidores centrais.
	
	Além disso, a educação dos utilizadores sobre os riscos associados ao partilhar dados emocionais e a promoção de controlos sobre como as suas informações são usadas são cruciais para capacitar os indivíduos e fomentar interacções mais seguras.
	
	\subsection{Integração Ética de Heurísticas Emocionais na IA}
	
	\textcite{gigerenzer2009} sugerem que as heurísticas humanas evoluíram para lidar com incertezas. A integração de heurísticas emocionais em sistemas de IA apresenta oportunidades e riscos. Para integrar heurísticas emocionais de forma ética, é necessário:
	
	\begin{itemize}
		\item \textbf{Transparência nos algoritmos}: Desenvolver algoritmos cujos processos de tomada de decisão possam ser auditados e compreendidos.
		\item \textbf{Diversidade nos dados de treino}: Utilizar conjuntos de dados que representem a diversidade cultural e emocional da população.
		\item \textbf{Supervisão humana}: Manter supervisão nas decisões críticas.
		\item \textbf{Ética incorporada}: Incluir princípios éticos no design dos sistemas, alinhados com normas internacionais como as propostas pelo \textcite{ieee2019}.
		\item \textbf{\textit{Feedback} contínuo}: Implementar mecanismos para que os utilizadores forneçam \textit{feedback}.
		\item \textbf{Avaliação contínua}: Realizar avaliações periódicas dos sistemas para garantir alinhamento com valores éticos.
	\end{itemize}
	
	\section{Conclusão}
	
	As emoções e a linguagem desempenham papéis cruciais na cognição humana e na tomada de decisão através de heurísticas. Embora possam facilitar decisões rápidas e adaptativas, também introduzem desafios significativos para a criação de sistemas de IA emocional. A complexidade e ambiguidade da linguagem humana representam obstáculos técnicos substanciais, enquanto questões éticas e de privacidade exigem atenção cuidadosa.
	
	Para maximizar os benefícios e mitigar os riscos, é essencial que desenvolvedores, reguladores e utilizadores colaborem na formulação de soluções éticas e transparentes. Conforme enfatizado por \textcite{russell2020}, o futuro da IA emocional depende da sua capacidade de respeitar os direitos humanos e promover interacções mais empáticas e seguras.
	
	O progresso nesta área requer um compromisso contínuo com a investigação interdisciplinar. Somente através de uma abordagem holística poderemos desenvolver sistemas de IA emocional que beneficiem a sociedade como um todo, respeitando a dignidade e a privacidade dos indivíduos. A inovação responsável nesta área tem o potencial de transformar positivamente a forma como interagimos com a tecnologia.
	
	Em última análise, a integração bem-sucedida de emoções e heurísticas na IA dependerá da nossa capacidade de compreender profundamente a natureza humana e de traduzir esse conhecimento em sistemas tecnológicos que reflitam não apenas a inteligência, mas também a sabedoria e a ética humanas.
	
	\newpage
	
	% Referências
	\printbibliography
	
\end{document}
